% from ega4-11-fr.tex
\oldpage[IV-3]{138}

les flèches verticales sont des isomorphismes~; comme la première ligne est un
isomorphisme, il en est de même de la seconde, donc la suite
$((g_i\otimes1)_x)$ est $(\mathcal F_y)_x$-quasi-régulière, et par suite aussi
$(\mathcal F_y)_x$-régulière, puisque $X_y$ est localement de type fini sur
$k(y)$. C.Q.F.D.
\end{proof}

\begin{theorem}[11.3.10]
\EGAReaderTarget{ega4.c04.thm.11-3-10}{IV.11.3.10-fr}
(critère de platitude par fibres). --- Soient $S$ un préschéma,
$g:X\to S$, $h:Y\to S$ deux morphismes, $f:X\to Y$ un $S$-morphisme,
$\mathcal F$ un $\mathcal O_X$-Module quasi-cohérent, $x$ un point de $X$,
$y=f(x)$, $s=h(y)=g(x)$. On suppose vérifiée l'une des deux hypothèses
suivantes~:
\begin{enumerate}[label=\arabic{enumi}\textsuperscript{o}]
\item $S$, $X$ et $Y$ sont localement noethériens et $\mathcal F$ est cohérent.
\item $g$ et $h$ sont localement de présentation finie et $\mathcal F$ est de
présentation finie.
\end{enumerate}
Alors, avec les notations de (9.4.1), si $\mathcal F_x\neq0$, les conditions
suivantes sont équivalentes~:
\begin{enumerate}[label=\emph{\alph*)}]
\item $\mathcal F$ est $g$-plat au point $x$ et $\mathcal F_s$ est $f_s$-plat
au point $x$.
\item $h$ est plat au point $y$ et $\mathcal F$ est $f$-plat au point $x$.
\end{enumerate}
En outre, dans l'hypothèse 2\textsuperscript{o}, l'ensemble des $x\in X$
vérifiant la condition b) est ouvert dans $X$.
\end{theorem}

\begin{proof}
La dernière assertion résulte de (11.3.1) appliqué à $\mathcal O_Y$ et au
morphisme $h$ d'une part, à $\mathcal F$ et au morphisme $f$ (qui est
localement de présentation finie) de l'autre.

Comme $g=h\circ f$, il est clair que b) implique a) sans supposer
1\textsuperscript{o} ni 2\textsuperscript{o} (2.1.6 et 2.1.4). Pour prouver que
a) entraîne b), on peut se borner au cas où $S$, $X$ et $Y$ sont affines~;
sous l'hypothèse 2\textsuperscript{o}, en appliquant (11.2.7), on se ramène au
cas où de plus $S$ est noethérien, c'est-à-dire qu'on peut se borner à
considérer le cas où l'hypothèse 1\textsuperscript{o} est satisfaite. Alors
l'assertion équivaut au lemme suivant, qui améliore
$(\mathbf 0_{\mathrm{III}},10.2.5)$~:

\begin{lemma}[11.3.10.1]
\label{IV.11.3.10.1-fr}
Soient $\rho:A\to B$, $\sigma:B\to C$ des homomorphismes locaux d'anneaux
locaux noethériens, $k$ le corps résiduel de $A$, $M$ un $C$-module
$\neq0$ de type fini. Alors les conditions suivantes sont équivalentes~:
\begin{enumerate}[label=\emph{\alph*)}]
\item $M$ est un $A$-module plat et $M\otimes_Ak$ est un
$(B\otimes_Ak)$-module plat.
\item $B$ est un $A$-module plat et $M$ est un $B$-module plat.
\end{enumerate}
\end{lemma}

Nous établirons d'abord le lemme plus général suivant~:

\begin{lemma}[11.3.10.2]
\label{IV.11.3.10.2-fr}
Soient $A$ un anneau commutatif, $B$ une $A$-algèbre commutative,
$\mathfrak J$ un idéal de $A$, $M$ un $B$-module. On considère d'une part les
conditions suivantes~:
\begin{enumerate}[label=\emph{(\roman*)}]
\item $\mathfrak J$ est nilpotent.
\item $B$ est noethérien et $M$ est idéalement séparé pour la topologie
$\mathfrak JB$-préadiques $(\mathbf 0_{\mathrm{III}},10.2.1)$.
\item $\mathfrak JB$ est contenu dans le radical de $B$.
\end{enumerate}
On considère d'autre part les quatre propriétés~:
\begin{enumerate}[label=\emph{\alph*)}]
\item $M$ est un $B$-module plat.
\item $B$ est un $A$-module plat.
\item $M$ est un $A$-module plat et $M/\mathfrak JM$ un
$(B/\mathfrak JB)$-module plat.
\item $B/\mathfrak JB$ est un $(A/\mathfrak J)$-module plat et pour tout idéal
maximal $\mathfrak m\supset\mathfrak JB$ de $B$ on a $\mathfrak mM\neq M$.
\end{enumerate}
Alors~:
\begin{enumerate}[label=\arabic{enumi}\textsuperscript{o}]
\item Si l'une des conditions (i), (ii) est vérifiée, la conjonction de a) et
b) implique c), et c) implique a).

